\documentclass[11pt,a4paper]{article}
\usepackage[T1]{fontenc}
\usepackage[spanish]{babel}
\usepackage[margin=2.4cm]{geometry}
\usepackage{tabularx,array}
\usepackage{graphicx}
\usepackage[hidelinks]{hyperref}
\usepackage[backend=biber,style=ieee]{biblatex}
\addbibresource{referencias.bib}
\setlength{\parindent}{0pt}
\setlength{\parskip}{5pt}
\renewcommand{\arraystretch}{1.25}
\newcommand{\pendiente}[1]{\textit{[#1]}}
\newcolumntype{Y}{>{\raggedright\arraybackslash}X}
\begin{document}
\begin{center}
\vspace*{.15cm}
{\large UNIVERSIDAD TÉCNICA ESTATAL QUEVEDO\par}
\vspace{.7cm}
Ingeniería de Software Aplicaciones Distribuidas\par
\vspace{.65cm}
Diseño de microservicios de\par
\vspace{.65cm}
Sistema de Gestión de Reservas de Laboratorio\par
\vspace{.65cm}
TOAQUIZA ZAMBRANO BELINDA BETZABETH\par
\vspace{.65cm}
PIGUAVE MONTAÑO NIXON ANTONIO\par
\vspace{.65cm}
EMANUEL PINO JULIANA ROMINA\par
\vspace{.65cm}
HERRERA BARCO HUMBERTO ALDAIR\par
\vspace{.65cm}
GA-TL-02 / FORM\par
\vspace{.65cm}
Grupo EHPT\par
\vspace{.65cm}
Vista previa de la estructura\par
\vspace{.65cm}
REPOSITORIO\par
\vspace{.65cm}
\href{https://github.com/JulianaEmanuel/GA-TL-02-FORM-Taller.git}{https://github.com/JulianaEmanuel/GA-TL-02-FORM-Taller.git}
\end{center}
\thispagestyle{empty}
\newpage
\section{Introducción}
El taller documenta el diseño de componentes de CLAB a partir de la propuesta de arquitectura de la entrega 1.

\section{Diseño de los servicios}
\subsection{Reservas y Agenda \hfill Juliana Emanuel}
\pendiente{Describir los componentes de disponibilidad, reservas, bloqueos y asistencia, y los datos de agenda que conserva el servicio.}
\subsection{Identidad y Académico \hfill Humberto Herrera}
\pendiente{Describir los componentes de acceso, permisos y gestión académica, y los datos propios del servicio.}
\subsection{Laboratorios y Recursos \hfill Belinda Toaquiza}
\pendiente{Describir los componentes de catálogo, inventario y seguimiento de fallas, y los datos propios del servicio.}
\subsection{Información y Avisos \hfill Nixon Piguave}
\pendiente{Describir los componentes de reportes, avisos y consulta central de auditoría, y los datos que mantiene el servicio.}

\section{API Gateway y comunicación}
\pendiente{Humberto documenta el acceso web y móvil, el enrutamiento, la autenticación y los límites de solicitudes. Cada responsable identifica sus canales, indica si son síncronos o asíncronos y justifica la tecnología según los requisitos no funcionales.}

\section{Diagrama C4 de componentes}
\IfFileExists{c4-nivel-3.pdf}{%
  \includegraphics[width=\linewidth,height=.55\textheight,keepaspectratio]{c4-nivel-3.pdf}%
}{\pendiente{Incorporar la exportación del diagrama integrado y registrar la ubicación del archivo editable.}}

\newpage
\section{Tabla de servicios}
La asignación inicial de requisitos sigue las historias de usuario de E1. Cada responsable completa la propiedad de datos y el contrato de su servicio.

\noindent\begin{tabularx}{\linewidth}{|p{3.3cm}|Y|Y|Y|}
\hline
\textbf{Servicio y responsable} & \textbf{Responsabilidad y datos propios} & \textbf{Endpoints o eventos} & \textbf{Requisitos funcionales}\\ \hline
Reservas y Agenda (Juliana) & \pendiente{Completar} & \pendiente{Completar} & RF-04 a RF-06, RF-11 y RF-12\\ \hline
Identidad y Académico (Humberto) & \pendiente{Completar} & \pendiente{Completar} & RF-01, RF-02 y RF-16\\ \hline
Laboratorios y Recursos (Belinda) & \pendiente{Completar} & \pendiente{Completar} & RF-03 y RF-07 a RF-10\\ \hline
Información y Avisos (Nixon) & \pendiente{Completar} & \pendiente{Completar} & RF-13, RF-15 y vista central de RF-14\\ \hline
\end{tabularx}

\textbf{RF-14.} Cada servicio registra sus cambios. Información y Avisos reúne la vista central.

\section{Retroalimentación de la revisión}
\noindent\begin{tabularx}{\linewidth}{|Y|Y|Y|}
\hline
\textbf{Observación} & \textbf{Cambio realizado en el diagrama} & \textbf{Razón si no se aplicó}\\ \hline
\pendiente{Registrar tras la revisión} & & \\ \hline
\end{tabularx}

\section{Evidencias}
\pendiente{Añadir la foto fechada del bosquejo grupal. Enlazar los ADR si la revisión cambia una decisión de arquitectura.}

\section{Declaración de IA}
Se usó OpenAI Codex para preparar la estructura de esta plantilla. \pendiente{Completar la versión y la verificación del equipo antes de entregar.}

\printbibliography[title={Referencias}]
\end{document}